\section{如何理解评估}

这节课给出的基本框架可以压缩成四个问题：

\begin{enumerate}
    \item 输入是什么？
    \item 如何调用模型？
    \item 如何评估输出？
    \item 如何解释结果？
\end{enumerate}

\subsection{目标先于指标}

同一个模型分数，可能对应完全不同的评价目标。

\begin{itemize}
    \item \textbf{用户或公司做购买决策}：该选 Claude、Gemini、Grok，还是别的模型？
    \item \textbf{研究者衡量能力}：模型是否在推进``原始智能''或语言能力？
    \item \textbf{政策和业务分析}：模型带来的收益和危害分别是什么？
    \item \textbf{模型开发者调参}：某个改动是否真的让模型更好？
\end{itemize}

\begin{warningbox}{为什么``我只是评估一下模型''这句话不够}
如果不先说清楚目标，评估结果通常会把不同东西混在一起。一个 benchmark 分数可能同时受到 prompt、解码策略、工具使用、训练数据污染、参考答案质量和后处理流程的影响。
\end{warningbox}

\subsection{输入、调用、输出、解释}

\textbf{输入}部分要问的是：这个评估覆盖了哪些使用场景？有没有尾部难例？输入是不是被改造成更适合模型的形式，比如多轮对话？

\textbf{调用}部分要问的是：是 zero-shot、few-shot，还是 chain-of-thought？有没有工具、RAG 或者其他 agent scaffolding？你评估的是裸模型，还是整个系统？

\textbf{输出}部分要问的是：参考答案是否干净？用什么 metric？像代码任务常用 pass@k，开放式生成可能需要人工偏好、LLM-as-judge 或别的办法。成本要不要算进去？医疗、法律这类场景里，错误是不是有不同代价？

\textbf{解释}部分要问的是：一个 91\% 的分数到底意味着什么？能直接上线吗？测试集和训练集有没有重叠？你在评估最终模型，还是某种训练方法？

\begin{importantbox}{多轮对话和红队测试会改变输入定义}
对于 chat 系统或 red teaming，输入本身往往需要针对模型动态调整。这样更贴近真实行为，但也让不同模型之间的横向对比更难，必须明确规则。
\end{importantbox}

